\begin{example}
\label{ex:ch25-fixed-reference-small-capacities}
There is no positive lower bound for the capacities of all divisorial
points even when the curve and reference form are fixed.  Take
\(C=\mathbf P^1_K\) and let \(\theta\) be the form of the standard
model \(\cO_{\mathbf P^1_R}(2)\).  Its measure is
\(\mu_0=2\delta_g\), where \(g\) is the Gauss point.  We keep this same
form throughout the construction.

Let \(T\) be the affine coordinate.  Blow up the point \(T=0\) on the
special fibre, then repeatedly blow up the intersection of the strict
transform of the section \(T=0\) with the newest exceptional component.
Each centre is a smooth residue-rational point of the special fibre, not
a node.  After \(N\geq1\) blowups, the regular model has a chain of
\(N+1\) multiplicity-one components, joined by \(N\) residue-rational
nodes with equation \(xy=\pi\).  Thus every edge has length and weight
one, and its skeleton is \(\Gamma_N=[0,N]\), with \(g\) at zero.
The newest endpoint \(v_N\) is divisorial; in the local coordinates of
the blowups its normalized valuation satisfies \(v_N(T)=N\).

The pullback of the fixed reference bundle has degree two on the strict
transform of the original component and degree zero on every exceptional
component, by the projection formula.  Consequently the endpoint masses
are \(m_0=2\), \(m_1=0\), and the interior masses vanish.  If \(p_N\)
is the retraction to \(\Gamma_N\), the endpoint formula in
\Cref{prop:ch25-segment-capacities}, applied with the orientation
reversed, gives
\[
 \Capa_\theta\bigl(p_N^{-1}(\{v_N\})\bigr)
                 =\min\{2,1/N\}=1/N.
\]
For completeness the lower bound is attained by the rational model
potential whose skeletal restriction is \(f_N(s)=-s/N\).  On this
regular model its vertex values \(-j/N\), for \(0\leq j\leq N\),
are represented by a vertical rational Cartier divisor.  The associated
potential is \(f_N\circ p_N\), lies in \([-1,0]\), and has measure
\[
 \MA_\theta(f_N\circ p_N)
       =(2-1/N)\delta_g+(1/N)\delta_{v_N}.
\]
Both coefficients are nonnegative, so this is an admissible
\(\theta\)-psh model potential.  Its mass at \(v_N\), together with
monotonicity and the displayed tube upper bound, proves
\[
          \Capa_\theta(\{v_N\})=1/N\longrightarrow0.
\]
The varying objects are the divisorial points and their determining
models, not \(C\), \(\theta\), or the capacity under consideration.
\end{example}
